\section{Concrete Lean integration, with no new global controller}
\label{sec:integration}

\subsection{Reuse Forge's boundary; add worker-specific semantics}

Forge already specifies a neutral intermediate representation, exact origin
ownership, worker isolation, and checked-fact publication. Repeating those
requirements at length would not improve the design~\cite{forge-integration,forge-transfer}.
The new implementation work is the mathematical meaning of the three prepared
problem types and the proofs connecting them to source expressions.

\begin{center}
\begin{tikzpicture}[node distance=7mm and 7mm,box/.style={draw=ink,rounded corners=2pt,align=center,font=\small,minimum height=10mm,text width=36mm},>=Stealth]
\node[box] (source) {Lean source goal\\and exact context};
\node[box,right=of source] (model) {Prepared model\\plus denotation proof};
\node[box,right=of model] (check) {Lean certificate\\checker and theorem};
\node[box,below=of model] (search) {Untrusted exact search\\Python / CAS};
\draw[->] (source)--(model);
\draw[->] (model)--(check);
\draw[->] (model)--(search);
\draw[->] (search.east)-|(check.south);
\end{tikzpicture}
\end{center}

The denotation proof belongs to preparation, before the external answer exists.
The external search supplies coefficients and indices, not executable Lean
source. Its positive result has to satisfy the checker's own reconstructed
identities for the already prepared model. The final proof is transported back
to the original goal. This instantiates Forge's existing boundary rather than
inventing a competing one.

\subsection{The smallest useful word-fold reifier}

Begin with equalities of rational-valued computations over the \emph{same}
\code{List Rat} input and fixed rational initial states. Recognise either an
explicit state fold or a pair of folds that can be synchronised into a product
state. The first implementation can restrict all state fields to rationals and
all updates to affine expressions in those fields. Do not start with arbitrary
recursive Lean functions.

A practical preparation algorithm is:

\begin{lstlisting}
prepareWordGoal(goal):
    expose the equality and its exact quantified list binder
    recognize approved fold definitions and their traversal direction
    synchronize same-input folds using a proved product-fold lemma
    reify initial state, update, and output difference
    reject unsupported operations and unresolved coefficient parameters
    choose the output-degree monomial lift
    generate update and observation denotation proofs
    return fixed matrices, variable origins, and correspondence theorem
\end{lstlisting}

The expression reifier accepts rational literals, addition, subtraction,
multiplication, and fixed natural powers, with degree bookkeeping. Division is
initially accepted only when it is part of a rational literal with a proved
nonzero numeral denominator. A multiplication of two state-dependent affine
expressions fails the affine-update test, although it can be valid in the
observer. A polynomial input variable can occur in update coefficients; its
post-lift degree must be calculated, not guessed.

A fold-right computation and a fold-left computation are not synchronised by
merely using the same list variable. Their relationship needs a proved
traversal lemma. Likewise a mapped, reversed, filtered, or indexed input is not
the original word unless a source-level transformation theorem justifies that
view. The initial source subset should refuse these cases, then add them with
explicit correspondence lemmas.

For a goal with arbitrary initial parameters, there are two principled future
routes. One can introduce a finite polynomial coefficient family of initial
rows and share a reachable-space closure across those rows. Or one can ask
Forge's existing generalisation planner to propose and prove a stronger state
invariant. The current runner uses concrete initial states; it does not
implement either extension as a source tactic.

\subsection{The matrix soundness theorem is smaller than the search}

In Lean, represent the rectangular certificate matrix using finite index types,
for example \code{Matrix (Fin b) (Fin d) Rat}. Define the embedding
\[
\operatorname{embed}(x)=xB,
\]
the full transition $s\mapsto sM_a$, and the small transition $x\mapsto xC_a$.
The identity $BM_a=C_aB$ proves that the embedding commutes with transition.
The identity $Bv=0$ proves the observation vanishes on the embedded state.
This is exactly the generic word-simulation argument.

The bundle includes \code{lean/WordSimulation.lean}, a proposed core Lean proof
of that generic argument, with no matrix search hidden in it. It is marked
\emph{uncompiled}; the working environment did not provide Lean. It is not a
reflected checker, a tactic, or an evidence receipt. Its intended use is to
separate the straightforward induction lemma from the subsequently required
matrix implementation.

For a first compiled implementation, construct a proof of each finite matrix
identity with ordinary rational normalisation and instantiate the soundness
lemma. For larger certificates, implement a reflected Boolean checker and
prove its correctness once. The latter optimisation should come after a real
small-certificate end-to-end path works. Exact elimination need never be
verified inside Lean merely to check the returned matrix identities.

\subsection{Summation preparation and residual obligations}

A summation preparer must recognise the exact range, integer exponent $p$,
rational weight $q$, cast structure, and binomial family. The first entry point
should accept the literal form \eqref{eq:sum}; equivalent-looking factorial
quotients become eligible only through proved rewrite lemmas. It returns the
family parameters and a theorem equating the source sum with its prepared
semantics.

The checked telescope publishes a recurrence theorem as an intermediate fact.
It should not pretend to close an arbitrary requested closed form. Forge can
then use an available target recurrence, ask an appropriate library adapter for
one, or leave that target-side recurrence as an explicit obligation. If the
source goal already asks for the recurrence, publication closes that part
directly.

The recurrence equality worker should generate a proof graph containing both
annihilation proofs, the operator-product equality proof, the singularity cover,
and every seed equality. A seed proof may come from \code{norm_num}, direct
finite computation through the kernel, a library lemma, or another worker.
If one seed remains unresolved, the result is a partial proof plan, not a
positive theorem result. This is a direct use of Forge's existing obligation
graph, not a new graph data structure.

\subsection{Proposed modules and concrete responsibilities}

The following paths are proposed additions to Forge; they do not denote files
already installed in the repository.

\begin{longtable}{@{}p{0.40\linewidth}p{0.55\linewidth}@{}}
\toprule
Module & Responsibility and acceptance invariant\\
\midrule\endhead
\path{Forge/Summary/WordModel.lean} & Rational row-state semantics; word evaluation; finite alphabet identity.\\
\path{Forge/Summary/WordCertificate.lean} & Typed $B,\alpha,C$ data; dimension-safe identity checker; proved soundness.\\
\path{Forge/Summary/PolynomialInput.lean} & Coefficient-matrix semantics; interpolation theorem; actual-grid counterexample reconstruction.\\
\path{Forge/Reify/AffineFold.lean} & Exact approved fold grammar; monomial lift; source denotation proof; explicit refusals.\\
\path{Forge/Summary/BinomialSupport.lean} & Signed-lower-index binomial; support, ratio, complement, and endpoint lemmas.\\
\path{Forge/Summary/Telescoper.lean} & Polynomial identities and boundary partition; recurrence theorem for the fixed sum family.\\
\path{Forge/Summary/ShiftOperator.lean} & Polynomial Ore product and sequence action; left-multiple transport theorem.\\
\path{Forge/Summary/RecurrenceSeeds.lean} & Complete natural-singularity cover and strong-induction equality theorem.\\
\path{Forge/Worker/FiniteSummary.lean} & Bind prepared goals to solver requests; check results; publish proofs or residual obligations through existing interfaces.\\
\bottomrule
\end{longtable}

Keep the first three mathematical modules independent of \grind{} internals.
Their output is an ordinary theorem. A controller can feed that theorem into a
fresh local \grind{} call through Forge's already proposed conservative bridge.
No changes to congruence closure or the Boolean search loop are necessary for
these first capabilities.

\subsection{Library map: use existing mathematics and search engines}

\begin{longtable}{@{}p{0.25\linewidth}p{0.38\linewidth}p{0.30\linewidth}@{}}
\toprule
Library & Concrete use & Status in this study\\
\midrule\endhead
Lean core / Mathlib & Finite indices, lists, matrices, rational arithmetic, polynomial evaluation, finite sums, choose identities, and proof tactics & API/module documentation inspected; no local Lean build.\\
SymPy 1.14.0 & Exact polynomial expansion and rational nullspaces for bounded telescoping and operator search & Executed by the delivered producer.\\
Python standard library & Rational matrix and sparse-polynomial certificate checking with \code{Fraction}; canonical JSON decoding & Executed in isolated \code{python -S} replay.\\
\code{python-flint} & Optional fast dense rational matrices, inversion, characteristic polynomials, and elimination backend & Official \code{fmpq_mat} documentation consulted; adapter not implemented.\\
Sage \code{ore_algebra} & Optional LCLM, operator arithmetic, and desingularisation search, returning polynomial certificates & Official capabilities inspected; not installed or benchmarked here.\\
HolonomicFunctions & Optional broader multivariate creative-telescoping search and Ore computations & Official package description inspected; no Wolfram execution in this study.\\
\bottomrule
\end{longtable}

The referenced Mathlib matrix-to-linear-map and span modules contain the
infrastructure relevant to the simulation theorem. The choose-sum module
already includes direct binomial identities, so those should be tried before
expensive rediscovery~\cite{mathlib-matrix,mathlib-span,mathlib-choose}.
SymPy's matrix API supplies the nullspace operation used here; its role is
proposal generation, not proof acceptance~\cite{sympy}.
The rational matrix type in \code{python-flint} is an appropriate candidate for
larger exact searches, but no speedup has been measured~\cite{flint}.

Forge's Lean-status ledger pins Lean 4.34.0 and a Mathlib revision, while also
warning that other proposed dependencies are not automatically build-compatible
with that pin~\cite{forge-lean}. The Lean release itself is recorded as
4.34.0~\cite{lean-release}. The new workers should use the repository's resolved
project environment rather than independently chasing current versions of all
libraries. The delivered Python requirement is separately pinned to the actual
executed version. No fabricated Lake manifest is supplied.

\section{A specific Leant/Djex extension: quantified equivalence of eligible folds}
\label{sec:leant}

\subsection{The gap to target}

Leant already generates source-typed candidates and checks supplied behavioral
assertions on the exact candidate. Its documented assertion strategy uses
positive and negative decision attempts and bounded simplification; finite
observations establish only the stated observations. A universal assertion can
remain inconclusive when it needs a mathematical argument beyond those
procedures~\cite{leant-behavior}.
Djex's source-typed evidence graph preserves exact source ownership and the
relationship between a candidate and its checked graph~\cite{djex-graph}.
Forge already incorporates the generic ownership lessons from those
projects~\cite{forge-transfer}. The addition proposed here is a particular
semantic consumer: an exact affine-fold equivalence check.

Given two eligible rational-valued list folds at the same source type, translate
each to a linear word model with correspondence proofs. Form their block
comparison model and invoke Worker I. A valid closure certificate gives equality
for every list, not just the finite list test set. A separating word is a useful
new behavioral discriminator, subject to the source correspondence discipline.

\subsection{An operational candidate pipeline}

\begin{lstlisting}
for each exact, source-typed candidate occurrence:
    try the proved affine-fold interpretation
    if unsupported: preserve the ordinary behavioral-check route
    otherwise compare to the approved specification model
    if a matrix closure certificate is returned:
        kernel-check it and transport to the exact candidate proposition
    if a separating word is returned:
        reconstruct the actual source input and verify its observed failure
    preserve candidate ownership, original budgets, and residual goals
\end{lstlisting}

The specification model must not be inferred from the candidate's own type
alone. A type such as \code{List Rat -> Rat} admits many different behaviors.
The user assertion or an approved specification theorem determines the intended
observation. Likewise, a candidate's exact occurrence must not borrow a graph
or proof receipt from another candidate with the same rendering.

The source bridge should begin with ordinary list folds, not arbitrary
polymorphic Church encodings. A theorem about encoded finite lists is not
automatically a theorem about all inhabitants of an abstract Church type.
Extending the bridge to those inputs needs the appropriate source semantics or
parametricity argument; it must not be assumed from a type-directed search
result. The original Leant integer-indexing and simultaneous-universe search
queries are outside the tested fragment and are not claimed solved here.

\subsection{Semantic deduplication that does not discard evidence}

A proved functional equality can justify sharing future \emph{extensional}
behavioral checks between eligible candidates. Maintain an equivalence class
whose edges are actual equalities, while retaining each candidate's own origin
and source typing evidence. A chosen representative does not acquire the
others' metadata. Transport an assertion along an equality only when the
assertion has the correct source type and Lean accepts that transport.

It is unsafe to use a matrix-model equality as a universal deduplication key for
arbitrary typed programs. The translation may cover only one observer, one
input domain, or one source specialisation. Even a full extensional equality is
not definitional equality and does not justify rewriting unrelated provider or
universe metadata. The useful optimization is narrower: avoid repeating a
behavioral proof that can actually be transported through a checked equality.

\subsection{Counterexamples improve search without becoming oracles of truth}

For an exact model, a separating word can be returned directly with a proof of
the relevant source inequality once the source bridge is in place. For an
over-approximate model, rerun the candidate on the reconstructed input before
using the observation to reject it. An infeasible abstract transition sequence
must not eliminate a source-correct candidate.

This provides a concrete feedback loop: previously found counterexample words
are cheap tests for later candidates; surviving candidates still need the
universal matrix certificate. Passing the accumulated examples never replaces
that proof. The new work is the finite-dimensional universal argument, not the
already proposed general counterexample-guided synthesis architecture.
